% =====================================================================
\chapter{Обновление прошивки и SD-карты}
\label{ch:firmware}
% =====================================================================
\chapterintro
  {какие программы живут в пульте и в модели, как сделать резервную копию, как обновить
   EdgeTX, содержимое SD-карты и радиомодуль ExpressLRS}
  {компьютер с браузером Chrome или Edge, кабель USB-C \emph{с передачей данных}
   (некоторые кабели умеют только заряжать!)}

\section{Сколько «мозгов» в пульте и модели}

\begin{simple}
В твоей системе не одна программа, а несколько. Каждая живёт в своём «мозге» (микросхеме)
и обновляется по-своему. Главное правило: \textbf{радиомодуль в пульте и приёмник в модели
должны иметь совместимые версии}, иначе они не договорятся.
\end{simple}

\begin{figure}[H]
\centering
\begin{tikzpicture}[every node/.style={font=\sffamily\scriptsize}]
  % пульт
  \draw[etx,dashed,rounded corners=3pt] (-1.2,-1.6) rectangle (6.6,1.9);
  \node[font=\sffamily\small\bfseries,text=etx,anchor=west] at (-1.1,1.6) {в пульте};
  \pic[transform shape] at (0.4,0.3) {chip={STM32}};
  \node[align=center] at (0.4,-0.9) {\textbf{EdgeTX}\\меню, миксеры};
  \pic[transform shape,scale=1.1] at (2.9,0.3) {sdcard};
  \node[align=center] at (2.9,-0.9) {\textbf{SD-карта}\\модели, звуки,\\скрипты};
  \pic[transform shape] at (5.3,0.3) {chip={ESP32}};
  \node[align=center] at (5.3,-0.9) {\textbf{модуль ELRS}\\радиопередатчик};
  % модель
  \draw[okgreen!60!black,dashed,rounded corners=3pt] (8.0,-1.6) rectangle (13.4,1.9);
  \node[font=\sffamily\small\bfseries,text=okgreen!50!black,anchor=west] at (8.1,1.6) {в модели};
  \pic[transform shape,scale=1.3] at (9.4,0.3) {receiver};
  \node[align=center] at (9.4,-0.9) {\textbf{приёмник ELRS}};
  \pic[transform shape,scale=1.2] at (12.2,0.3) {fcboard};
  \node[align=center] at (12.2,-0.9) {\textbf{полётный}\\\textbf{контроллер}\\(Betaflight…)};
  \radiowaves{(6.6,0.3)}{0}
  \draw[darr,accent] (5.3,1.3) to[bend left=30] node[above,font=\sffamily\scriptsize\bfseries,text=accent,fill=white,inner sep=1pt]{версии должны совпадать} (9.4,1.3);
\end{tikzpicture}
\caption{Программы, которые может понадобиться обновить}
\end{figure}

\begin{center}\small
\renewcommand{\arraystretch}{1.25}
\begin{tabularx}{\linewidth}{@{}L{32mm}L{42mm}Y@{}}
\toprule
\textbf{Что} & \textbf{Чем обновлять} & \textbf{Когда} \\
\midrule
EdgeTX (пульт) & EdgeTX Buddy или загрузчик & нужна новая функция или исправление \\
SD-карта & EdgeTX Buddy или архив с сайта & \emph{всегда} после обновления EdgeTX \\
Модуль ELRS & ExpressLRS Configurator, по Wi-Fi & когда обновляешь приёмники \\
Приёмник ELRS & ExpressLRS Configurator, по Wi-Fi & вместе с модулем \\
Полётный контроллер & Betaflight / INAV Configurator & по необходимости \\
\bottomrule
\end{tabularx}
\end{center}

\section{Нужно ли вообще обновляться?}

Пульт приходит с рабочей версией --- летать можно сразу. Обновляйся, если:
\begin{itemize}
  \item в новой версии есть то, что тебе нужно;
  \item у модуля ELRS и приёмников разные версии;
  \item в группе или кружке все пульты должны быть одинаковыми.
\end{itemize}

\begin{caution}
Не обновляйся накануне соревнований или важного полёта. И не ставь на основной пульт
«ночные» (\emph{nightly}) версии --- это черновики для тестировщиков.
\end{caution}

\section{Шаг 0. Узнай версию}

\mpath{SYS\sep About} --- последняя страница меню пульта:

\begin{center}
\lcdscreen{ABOUT}{9/9}{\lblank\lc{EdgeTX}\lc{v2.11.2}\lc{tx12mk2}\lblank\lblank}
\end{center}

Строка \code{tx12mk2} --- это \emph{для какого пульта} собрана прошивка. У TX12, TX12 MKII и
Boxer прошивки \textbf{разные}: \code{tx12}, \code{tx12mk2}, \code{boxer}.

\section{Шаг 1. Резервная копия}

\begin{figure}[H]
\centering
\begin{tikzpicture}[every node/.style={font=\sffamily\small}]
  \pic[transform shape] at (0,0) {minitx};
  \draw[black!70,line width=1.4pt] (0,-0.55) .. controls (0,-1.4) and (2.5,-1.4) .. (2.8,-0.2);
  \node[lbl] at (1.3,-1.5) {USB-C};
  \pic[transform shape,scale=1.2] at (3.8,-0.4) {laptop};
  \node[font=\sffamily\scriptsize,align=center] at (3.8,0.35) {диск\\«EDGETX»};
  \draw[arr] (5.3,0.2) -- (6.6,0.2);
  \node[draw=black!50,fill=yellow!20,rounded corners=2pt,align=left,font=\ttfamily\scriptsize,anchor=west] at (6.8,0.2)
    {backup-2026-09-26/\\\ \ MODELS/\\\ \ RADIO/\\\ \ SOUNDS/ …};
\end{tikzpicture}
\caption{Копируем всю SD-карту в папку на компьютере}
\end{figure}

\begin{enumerate}
  \item Включи пульт и подключи к компьютеру. На экране пульта появится вопрос ---
        выбери \menu{USB Storage (SD)}.
  \item На компьютере появится новый диск. Скопируй \textbf{все} папки в новую папку с датой.
  \item Безопасно извлеки диск.
\end{enumerate}

\section{Шаг 2. Прошивка через EdgeTX Buddy}

\begin{figure}[H]
\centering
\begin{tikzpicture}[every node/.style={font=\sffamily\small}]
  \foreach \i/\t/\x in {1/{Открой\\buddy.edgetx.org\\в Chrome},2/{Выбери версию\\и свой пульт},3/{Выключи пульт\\и подключи USB},4/{Нажми «Flash»,\\жди 1--2 мин}}{
    \node[draw=etx,fill=etx!6,rounded corners=3pt,text width=30mm,minimum height=17mm,align=center] (b\i) at (\i*3.8,0) {\t};
    \node[circle,fill=accent,text=white,font=\sffamily\small\bfseries,inner sep=1.5pt] at ($(b\i.north west)+(0.15,-0.15)$) {\i};
  }
  \foreach \i [evaluate=\i as \j using int(\i+1)] in {1,2,3}{\draw[arr] (b\i) -- (b\j);}
\end{tikzpicture}
\caption{Обновление EdgeTX за 4 шага}
\end{figure}

\begin{enumerate}
  \item Открой \code{buddy.edgetx.org} в Chrome или Edge (другие браузеры не умеют работать с USB).
  \item Вкладка \menu{Flash}: выбери последнюю \emph{стабильную} версию (release) и пульт ---
        \menu{RadioMaster TX12 Mark II}, \menu{RadioMaster TX12} или \menu{RadioMaster Boxer}.
        Можно выбрать язык меню, в том числе русский.
  \item \textbf{Выключи} пульт и подключи его кабелем. Компьютер увидит устройство
        \code{STM32 BOOTLOADER} (экран пульта при этом может остаться тёмным --- это нормально).
  \item Нажми \menu{Flash}, выбери устройство в окошке браузера и жди. \textbf{Не выдёргивай кабель!}
\end{enumerate}

\begin{tip}
Если браузер не видит пульт на Windows, нужен драйвер WinUSB --- Buddy покажет, как его
поставить (обычно утилитой Zadig). Если кабель «только для зарядки», компьютер не увидит
вообще ничего --- возьми другой кабель.
\end{tip}

\section{Шаг 3. Обнови SD-карту}

Новой версии EdgeTX нужны новые звуки и скрипты. В Buddy вкладка \menu{SD Card}:
\begin{enumerate}
  \item включи пульт в режиме \menu{USB Storage (SD)};
  \item выбери диск пульта, версию и язык голоса (например, русский);
  \item Buddy допишет нужные файлы и \emph{не тронет} твои модели.
\end{enumerate}

Если забыть этот шаг, при включении пульт покажет предупреждение о версии SD-карты.

\section{Запасной путь: загрузчик}

Если Buddy не видит пульт, можно прошить со SD-карты:

\begin{figure}[H]
\centering
\begin{tikzpicture}[scale=0.5]
  \drawTXtwelve
  \foreach \x in {-4.9,4.9}{\path[fill=accent,draw=black,rounded corners=1pt] (\x-1.1,-1.55) rectangle (\x+1.1,-1.3);}
  \draw[motion] (-8.3,-1.42) -- (-6.1,-1.42);
  \draw[motion] (8.3,-1.42) -- (6.1,-1.42);
  \path[fill=accent] (0,-2.1) circle (0.3);
  \node[font=\sffamily\small\bfseries,text=accent,anchor=east] at (-8.4,-1.42) {1. прижми};
  \node[font=\sffamily\small\bfseries,text=accent,anchor=west] at (8.4,-1.42) {1. прижми};
  \draw[motion] (0,-4.3) -- (0,-2.5);
  \node[font=\sffamily\small\bfseries,text=accent,anchor=north] at (0,-4.3) {2. включи};
\end{tikzpicture}
\caption{Вход в загрузчик: прижми оба \emph{горизонтальных} триммера \emph{внутрь} (к центру)
и, не отпуская, включи пульт}
\end{figure}

\begin{enumerate}
  \item Скачай файл прошивки \code{.bin} для своего пульта и положи в папку \code{FIRMWARE}
        на SD-карте.
  \item Войди в загрузчик (рисунок выше).
  \item Выбери \menu{Write Firmware} → файл → подтверди. Дождись конца и выбери \menu{Exit}.
\end{enumerate}

\begin{center}
\lcdscreen{EdgeTX Bootloader}{}{\lblank\lsel{Write Firmware}{}\lr{Restore Firmware}{}\lr{Exit}{}\lblank\lblank}
\end{center}

\section{Обновление модуля ExpressLRS}
\label{sec:elrs-update}

\begin{figure}[H]
\centering
\begin{tikzpicture}[every node/.style={font=\sffamily\small}]
  \pic[transform shape,scale=1.3] at (0,0) {minitx};
  \node[lbl,align=center] at (0,-1.4) {пульт:\\ExpressLRS → WiFi};
  \foreach \r in {0.35,0.6,0.85}{\draw[etx,line width=1pt] (1.6,0.2) ++(-40:\r) arc (-40:40:\r);}
  \node[lbl] at (2.2,1.35) {Wi-Fi «ExpressLRS TX»};
  \pic[transform shape,scale=1.3] at (4.8,-0.3) {laptop};
  \node[font=\ttfamily\scriptsize] at (4.8,0.6) {10.0.0.1};
  \node[lbl,align=center] at (4.8,-1.4) {браузер: загрузи\\firmware.bin};
  \draw[arr] (7.0,0.2) -- (8.2,0.2);
  \node[draw=okgreen,fill=okgreen!8,rounded corners=2pt,align=center] at (9.6,0.2) {модуль\\обновлён};
\end{tikzpicture}
\caption{Обновление модуля ELRS по Wi-Fi}
\end{figure}

\begin{enumerate}
  \item Установи на компьютер \textbf{ExpressLRS Configurator}. Выбери версию и свой пульт
        (например, \menu{RadioMaster TX12 2400 TX} или \menu{RadioMaster Boxer 2400 TX}).
  \item Придумай \textbf{фразу привязки} (binding phrase) --- любые слова. Её же потом впишешь
        во все свои приёмники, и они будут связываться с пультом сами (глава~\ref{ch:rf}).
  \item Способ прошивки --- \menu{WiFi}. Нажми \menu{Build} --- получишь файл \code{firmware.bin}.
  \item На пульте: \mpath{SYS\sep Tools\sep ExpressLRS\sep WiFi Connectivity\sep Enable WiFi}.
  \item Подключи компьютер к Wi-Fi \code{ExpressLRS TX} (пароль \code{expresslrs}),
        открой \code{http://10.0.0.1} и загрузи файл.
\end{enumerate}

Есть и способ через USB: в Configurator выбери \menu{EdgeTX Passthrough}, а пульт подключи в
режиме \menu{USB Serial (VCP)}.

\subsection{Мультипротокольный модуль}

Для пультов с модулем «4-в-1»: скачай прошивку модуля (\code{.bin}), положи в папку
\code{FIRMWARE}, затем \mpath{SYS\sep SD card}, найди файл, долгое \btn{ENTER} →
\menu{Flash internal Multi}.

\begin{tryit}
\begin{enumerate}
  \item Сделай резервную копию SD-карты своего пульта.
  \item Посмотри версию EdgeTX в \menu{About} и сравни с последней на \code{edgetx.org}.
  \item Открой \mpath{SYS\sep Tools\sep ExpressLRS} и посмотри версию модуля
        (в нижней строке скрипта).
\end{enumerate}
\end{tryit}

\begin{summary}
\begin{itemize}
  \item Сначала --- резервная копия SD-карты.
  \item EdgeTX обновляется через Buddy (выключенный пульт + USB) или через загрузчик
        (триммеры внутрь + питание).
  \item После EdgeTX обязательно обнови SD-карту.
  \item Модуль ELRS и приёмники обновляются отдельно и должны иметь совместимые версии.
\end{itemize}
\end{summary}
